\documentclass[12pt,a4paper]{article}
\usepackage[utf8]{inputenc}
\usepackage[french]{babel}
\usepackage[T1]{fontenc}
\usepackage{geometry}
\usepackage{graphicx}
\usepackage{hyperref}
\usepackage{listings}
\usepackage{xcolor}
\usepackage{float}
\usepackage{amsmath}
\usepackage{titlesec}
\usepackage{fancyhdr}
\usepackage{tocloft}

% Configuration de la page
\geometry{margin=2.5cm}
\pagestyle{fancy}
\fancyhf{}
\fancyhead[L]{\leftmark}
\fancyhead[R]{\thepage}
\fancyfoot[C]{Rapport de Projet - Flint Commerce}

% Configuration des couleurs pour le code
\definecolor{codegreen}{rgb}{0,0.6,0}
\definecolor{codegray}{rgb}{0.5,0.5,0.5}
\definecolor{codepurple}{rgb}{0.58,0,0.82}
\definecolor{backcolour}{rgb}{0.95,0.95,0.92}

% Style pour les sections
\titleformat{\section}
{\Large\bfseries\color{blue!70!black}}
{\thesection}{1em}{}

\titleformat{\subsection}
{\large\bfseries\color{blue!50!black}}
{\thesubsection}{1em}{}

% Informations du document
\title{
    \vspace{-2cm}
    \textbf{\Huge Rapport de Projet}\\
    \vspace{0.5cm}
    \Large\textit{Flint Commerce - Plateforme E-Commerce}\\
    \vspace{0.3cm}
    \large Description Complète des Interfaces Utilisateur
}
\author{}
\date{\today}

\begin{document}

\maketitle
\thispagestyle{empty}

\newpage
\tableofcontents
\newpage

\section{Introduction}

\subsection{Contexte du Projet}

Flint Commerce est une plateforme e-commerce moderne et complète développée pour offrir une expérience d'achat en ligne optimale. Ce projet a été conçu avec une architecture full-stack utilisant les technologies web les plus récentes pour garantir à la fois des performances élevées et une expérience utilisateur exceptionnelle.

La plateforme se distingue par son interface utilisateur moderne et responsive, son système de gestion robuste pour les administrateurs, et ses fonctionnalités avancées de gestion de panier, de commandes et de paiement. L'objectif principal est de fournir une solution e-commerce complète qui répond aux besoins tant des clients que des administrateurs.

\subsection{Objectifs du Projet}

Les objectifs principaux de ce projet sont les suivants :

\begin{itemize}
    \item \textbf{Expérience Client Optimale} : Offrir une interface intuitive et moderne permettant aux clients de naviguer facilement, de rechercher des produits, de gérer leur panier et de passer des commandes en toute simplicité.
    
    \item \textbf{Gestion Administrative Complète} : Fournir aux administrateurs un tableau de bord complet avec des outils de gestion des produits, des commandes, des bannières et des analyses en temps réel.
    
    \item \textbf{Sécurité et Performance} : Assurer la sécurité des données utilisateur et des transactions, tout en maintenant des performances optimales grâce à une architecture bien structurée.
    
    \item \textbf{Responsive Design} : Garantir une expérience utilisateur cohérente et optimale sur tous les appareils (desktop, tablette, mobile).
\end{itemize}

\subsection{Architecture Technique}

Le projet est structuré en deux parties principales :

\begin{description}
    \item[\textbf{Frontend}] Application React.js développée avec Vite, utilisant Redux Toolkit pour la gestion d'état, React Router pour la navigation, et Tailwind CSS avec Shadcn/ui pour l'interface utilisateur moderne.
    
    \item[\textbf{Backend}] API REST développée avec Node.js et Express.js, utilisant MongoDB comme base de données via Mongoose, et intégrant Cloudinary pour la gestion des images.
\end{description}

\subsection{Structure du Document}

Ce rapport présente une description détaillée de toutes les interfaces utilisateur de la plateforme Flint Commerce. Chaque section décrit une interface spécifique, ses fonctionnalités principales, et son rôle dans l'écosystème global de l'application. Les descriptions sont accompagnées d'espaces réservés pour les captures d'écran qui illustreront visuellement chaque interface.

\newpage

\section{Architecture Générale du Système}

\subsection{Vue d'Ensemble}

Flint Commerce est organisé en trois espaces principaux distincts :

\begin{enumerate}
    \item \textbf{Espace Client} : Interface publique permettant aux utilisateurs de parcourir et d'acheter des produits.
    \item \textbf{Espace Administrateur} : Interface privée réservée aux administrateurs pour gérer la plateforme.
    \item \textbf{Espace Authentification} : Interface dédiée à la connexion et à l'inscription des utilisateurs.
\end{enumerate}

\subsection{Technologies Utilisées}

\subsubsection{Frontend}
\begin{itemize}
    \item React.js 18.3.1 - Bibliothèque JavaScript pour la construction d'interfaces utilisateur
    \item Vite 5.4.1 - Outil de build et serveur de développement
    \item Redux Toolkit 2.2.7 - Gestion d'état globale
    \item React Router DOM 6.26.2 - Routage côté client
    \item Tailwind CSS 3.4.11 - Framework CSS utility-first
    \item Shadcn/ui - Composants UI modernes basés sur Radix UI
    \item Lucide React - Bibliothèque d'icônes
    \item Recharts 2.13.3 - Bibliothèque de graphiques
\end{itemize}

\subsubsection{Backend}
\begin{itemize}
    \item Node.js - Environnement d'exécution JavaScript
    \item Express.js - Framework web pour Node.js
    \item MongoDB - Base de données NoSQL
    \item Mongoose - ODM (Object Data Modeling) pour MongoDB
    \item Cloudinary - Service de gestion d'images dans le cloud
    \item Nodemailer - Service d'envoi d'emails
\end{itemize}

\newpage

\section{Interfaces Client}

Cette section décrit toutes les interfaces accessibles aux clients de la plateforme. Ces interfaces sont conçues pour offrir une expérience d'achat fluide et intuitive.

\subsection{Page d'Accueil}

La page d'accueil (\texttt{/shop/home}) est la première interface que les visiteurs rencontrent. Elle sert de point d'entrée principal à la plateforme et présente les éléments essentiels pour attirer et guider les utilisateurs.

\subsubsection{Description}

La page d'accueil comprend plusieurs sections stratégiquement organisées pour offrir une expérience d'accueil complète et engageante. Un carousel automatique affiche les bannières promotionnelles configurées par les administrateurs, défilant automatiquement toutes les 5 secondes avec des contrôles de navigation manuels pour permettre aux utilisateurs de naviguer à leur rythme. La section Hero offre une accroche visuelle avec un design moderne utilisant des dégradés et des animations fluides pour créer un impact visuel fort dès l'arrivée sur la page.

La vitrine des marques présente une grille interactive montrant les différentes marques disponibles sur la plateforme. Chaque marque est représentée par une carte avec effet de survol, permettant aux utilisateurs de filtrer directement les produits par marque en un simple clic. Une section fonctionnalités met en avant trois avantages principaux de la plateforme : la garantie de qualité, les retours faciles, et le paiement sécurisé, rassurant ainsi les visiteurs sur la fiabilité de la plateforme.

La page d'accueil comprend également une section dédiée aux produits vedettes, mettant en avant les articles les plus importants ou les plus populaires, ainsi qu'une section affichant les derniers produits ajoutés au catalogue, permettant aux clients de découvrir les nouveautés en temps réel.

\subsubsection{Fonctionnalités}

La page d'accueil offre une navigation automatique du carousel avec possibilité de contrôle manuel, permettant aux utilisateurs de naviguer à leur convenance. Le filtrage direct par marque depuis la vitrine des marques facilite l'accès aux produits souhaités. L'accès rapide aux détails des produits via les cartes produits optimise le parcours utilisateur. Le design responsive est adapté à tous les écrans, garantissant une expérience optimale sur desktop, tablette et mobile. Les animations et transitions fluides contribuent à créer une expérience utilisateur agréable et moderne.

\subsubsection{Capture d'Écran}

\begin{figure}[H]
    \centering
    \fbox{\parbox{0.9\textwidth}{\centering\vspace{5cm}[Image de la page d'accueil]\vspace{5cm}}}
    \caption{Page d'accueil de Flint Commerce}
    \label{fig:home}
\end{figure}

\subsection{Page Catalogue}

La page catalogue (\texttt{/shop/catalog}) est l'interface principale de navigation et de recherche de produits. Elle permet aux utilisateurs de parcourir l'ensemble du catalogue avec des options de filtrage avancées.

\subsubsection{Description}

Cette page offre une expérience de navigation complète avec une barre latérale fixe (sticky) contenant tous les filtres disponibles. Les utilisateurs peuvent filtrer par catégorie (Hommes, Femmes, Enfants, Accessoires, Chaussures), par marque, par gamme de prix, par taille, et par couleur, permettant ainsi un affinage précis de la recherche. Une barre d'outils permet de trier les produits selon différents critères : prix (croissant ou décroissant), date d'ajout (plus récent ou plus ancien), popularité, ou nom (alphabétique).

La grille de produits affiche les résultats sous forme de cartes modernes et responsive. Chaque carte présente l'image du produit avec effet de survol, un badge de réduction si applicable avec le pourcentage affiché, le nom et la description du produit, ainsi que le prix avec le prix original barré si en promotion et le prix de vente mis en avant. Des badges "Nouveau" et "Vedette" permettent d'identifier rapidement les produits récents ou mis en avant.

Un système de pagination permet de naviguer entre les pages de résultats, et un message informatif et visuellement attrayant s'affiche lorsque aucun produit ne correspond aux filtres sélectionnés, guidant l'utilisateur vers une nouvelle recherche.

\subsubsection{Fonctionnalités}

La page catalogue offre un filtrage multi-critères avec application en temps réel, permettant aux utilisateurs de voir immédiatement les résultats de leurs sélections. Le tri dynamique des résultats s'adapte instantanément aux choix de l'utilisateur. La recherche visuelle rapide est facilitée grâce aux cartes produits claires et informatives. Les filtres sont persistants dans la session, permettant de revenir aux critères précédents facilement. Un compteur de résultats affiché en temps réel informe l'utilisateur du nombre de produits correspondants. Le design responsive avec sidebar rétractable sur mobile garantit une expérience optimale sur tous les appareils.

\subsubsection{Capture d'Écran}

\begin{figure}[H]
    \centering
    \fbox{\parbox{0.9\textwidth}{\centering\vspace{5cm}[Image de la page catalogue]\vspace{5cm}}}
    \caption{Page catalogue avec filtres et grille de produits}
    \label{fig:catalog}
\end{figure}

\subsection{Page Détails Produit}

La page de détails produit (\texttt{/shop/product/:id}) fournit toutes les informations nécessaires pour qu'un client puisse prendre une décision d'achat éclairée.

\subsubsection{Description}

Cette page présente le produit de manière exhaustive pour permettre au client de prendre une décision d'achat éclairée. La galerie d'images offre une image principale en grande taille avec possibilité de zoom, accompagnée de miniatures cliquables pour naviguer entre les différentes images du produit. Les transitions fluides entre les images et le support multi-images pour chaque produit permettent une visualisation complète de l'article.

Les informations produit sont présentées de manière claire et structurée. Le nom du produit est affiché en grand format, la marque et la catégorie sont indiquées avec des badges colorés, et une description détaillée informe le client sur les caractéristiques du produit. Le prix est affiché de manière proéminente, avec le prix original barré si le produit est en promotion, et le pourcentage de réduction est calculé et affiché automatiquement. Un indicateur visuel montre le stock disponible.

La sélection de taille utilise des boutons avec état visuel clair (sélectionné/non sélectionné), indique les tailles disponibles, et désactive automatiquement les tailles en rupture de stock. Les actions d'achat comprennent un bouton "Ajouter au Panier" avec animation de confirmation, un sélecteur de quantité avec boutons plus et moins, et une validation de la taille avant ajout au panier pour éviter les erreurs.

La section avis clients affiche les avis et notes des clients précédents, avec une note moyenne calculée et affichée avec des étoiles. Chaque avis présente le nom du client, la note attribuée et le commentaire. Les clients connectés peuvent ajouter leur propre avis, contribuant ainsi à la communauté. Enfin, une section produits similaires suggère des articles de la même catégorie ou marque, facilitant la découverte d'alternatives.

\subsubsection{Fonctionnalités}

La page de détails produit offre une navigation intuitive dans la galerie d'images avec zoom et navigation par miniatures. La validation de la sélection de taille avant ajout au panier évite les erreurs de commande. La gestion de la quantité respecte les limites de stock disponibles. L'affichage dynamique du prix s'adapte automatiquement selon les promotions en cours. Le système d'avis et de notation intégré permet aux clients de partager leur expérience et d'aider les autres acheteurs. Les suggestions de produits similaires facilitent la découverte d'articles complémentaires. Le design responsive est optimisé pour mobile, garantissant une expérience de qualité sur tous les appareils.

\subsubsection{Capture d'Écran}

\begin{figure}[H]
    \centering
    \fbox{\parbox{0.9\textwidth}{\centering\vspace{5cm}[Image de la page détails produit]\vspace{5cm}}}
    \caption{Page de détails produit avec galerie et informations}
    \label{fig:product-details}
\end{figure}

\subsection{Page Panier}

Le panier est accessible via un panneau latéral (sheet) qui s'ouvre depuis la barre de navigation. Il permet aux utilisateurs de gérer leurs articles avant de procéder au paiement.

\subsubsection{Description}

Le panier offre une interface compacte mais complète, accessible via un panneau latéral qui s'ouvre depuis la barre de navigation. L'en-tête présente le titre "Panier" avec un compteur d'articles et un bouton de fermeture pour une navigation intuitive. Chaque article est affiché sous forme de carte avec l'image miniature du produit, le nom et la taille sélectionnée, le prix unitaire et le prix total pour l'article, des contrôles de quantité avec boutons plus et moins, et un bouton de suppression avec confirmation pour éviter les suppressions accidentelles.

Le résumé du panier calcule automatiquement le sous-total, affiche les frais de livraison si applicables, et met en avant le total général de manière proéminente. Les actions disponibles incluent un bouton "Passer la Commande" menant directement à la page de paiement, et un bouton "Continuer les Achats" pour fermer le panier et retourner au catalogue.

Pour les utilisateurs non connectés, un message invite à se connecter ou s'inscrire, avec des boutons de redirection vers les pages d'authentification pour faciliter la création de compte. Lorsque le panier est vide, un message informatif suggère de parcourir le catalogue pour découvrir les produits disponibles.

\subsubsection{Fonctionnalités}

Le panier offre une mise à jour en temps réel des quantités et totaux, permettant aux utilisateurs de voir immédiatement l'impact de leurs modifications. Pour les utilisateurs non connectés, le panier est persisté dans le stockage local, permettant de conserver les articles même après fermeture du navigateur. Pour les utilisateurs connectés, la synchronisation avec le backend assure la cohérence des données sur tous les appareils. Le calcul automatique des totaux élimine les erreurs de calcul manuel. La validation avant suppression d'articles prévient les suppressions accidentelles. Le design responsive avec adaptation mobile garantit une expérience optimale sur tous les écrans.

\subsubsection{Capture d'Écran}

\begin{figure}[H]
    \centering
    \fbox{\parbox{0.9\textwidth}{\centering\vspace{5cm}[Image du panier]\vspace{5cm}}}
    \caption{Interface du panier d'achat}
    \label{fig:cart}
\end{figure}

\subsection{Page Paiement / Checkout}

La page de paiement (\texttt{/shop/checkout}) est l'interface finale du processus d'achat. Elle guide l'utilisateur à travers les étapes nécessaires pour finaliser sa commande.

\subsubsection{Description}

Cette page est organisée en un processus en plusieurs étapes avec un système visuel montrant la progression : l'étape 1 concerne les informations de livraison, et l'étape 2 le récapitulatif et paiement. La colonne principale est dédiée aux informations de livraison, permettant la sélection d'une adresse existante via une grille de cartes d'adresses. Chaque carte d'adresse présente un indicateur visuel de sélection avec bordure colorée, des icônes pour le type d'adresse (domicile, travail, autre), le nom complet du destinataire, l'adresse complète formatée, le numéro de téléphone, et des boutons d'édition et de suppression pour une gestion complète.

Un formulaire d'ajout et de modification d'adresse avec tous les champs nécessaires permet de créer ou modifier facilement les adresses de livraison. Un bouton "Ajouter une Adresse" avec style moderne facilite l'accès à cette fonctionnalité.

La colonne latérale présente le récapitulatif de commande en position sticky, restant visible lors du défilement. Elle affiche la liste des articles avec image miniature, nom du produit et taille, quantité, et prix unitaire et total. Les détails de prix incluent le sous-total, les frais de livraison, les réductions si applicables, et le total général affiché en grand format. Les informations de paiement montrent la méthode de paiement sélectionnée et le statut de validation. Le bouton "Passer la Commande" utilise un style avec dégradé moderne, affiche un état de chargement pendant le traitement, et valide l'adresse avant activation.

Le design moderne de la page utilise des fonds avec dégradés subtils pour créer une expérience visuelle agréable et professionnelle.

\subsubsection{Fonctionnalités}

La page de paiement offre un processus de commande en plusieurs étapes avec indicateurs visuels clairs pour guider l'utilisateur. La gestion complète des adresses permet l'ajout, la modification, la suppression et la sélection d'adresses de livraison. Le récapitulatif en temps réel calcule automatiquement les totaux, permettant à l'utilisateur de voir immédiatement l'impact de ses choix. La validation de l'adresse avant finalisation évite les erreurs de livraison. L'affichage sticky du récapitulatif garantit que les informations importantes restent toujours visibles. La gestion des erreurs avec messages clairs guide l'utilisateur en cas de problème. Après une commande réussie, la redirection automatique vers la page de commandes et le vidage automatique du panier complètent l'expérience utilisateur de manière fluide.

\subsubsection{Capture d'Écran}

\begin{figure}[H]
    \centering
    \fbox{\parbox{0.9\textwidth}{\centering\vspace{5cm}[Image de la page de paiement]\vspace{5cm}}}
    \caption{Page de paiement avec sélection d'adresse et récapitulatif}
    \label{fig:checkout}
\end{figure}

\subsection{Page Compte Utilisateur}

La page compte (\texttt{/shop/account}) permet aux utilisateurs de gérer leurs informations personnelles, consulter leurs commandes et gérer leurs adresses.

\subsubsection{Description}

Cette page est organisée en onglets pour une navigation claire et intuitive. La bannière de profil présente l'avatar utilisateur avec initiales ou image, le nom complet de l'utilisateur, l'email, le tout dans un design avec dégradé moderne qui crée une identité visuelle forte.

L'onglet Commandes affiche un tableau moderne listant toutes les commandes avec le numéro de commande formaté et cliquable, la date de commande formatée, le statut avec badge coloré (En attente, Confirmée, Expédiée, Livrée, Annulée), le montant total, les images empilées des produits commandés, et un bouton "Voir Détails" ouvrant un dialogue modal. Ce dialogue de détails présente les informations complètes de la commande, la liste détaillée des produits avec images, l'adresse de livraison, la méthode de paiement, et le récapitulatif des prix. Un état vide avec message s'affiche si aucune commande n'est présente.

L'onglet Adresses présente une grille de cartes d'adresses avec toutes les informations nécessaires. Chaque carte dispose de boutons d'édition et de suppression pour une gestion complète. Un formulaire d'ajout et de modification d'adresse avec validation des champs permet de maintenir à jour les informations de livraison.

\subsubsection{Fonctionnalités}

La page compte offre une navigation par onglets intuitive facilitant l'accès aux différentes sections. L'affichage complet de l'historique des commandes permet à l'utilisateur de suivre toutes ses transactions. Les détails de commande accessibles via dialogue modal offrent une vue détaillée sans quitter la page principale. La gestion complète des adresses (CRUD) permet un contrôle total sur les informations de livraison. Le design moderne avec cartes et badges crée une interface visuellement attrayante. Le responsive design pour mobile garantit une expérience optimale sur tous les appareils.

\subsubsection{Capture d'Écran}

\begin{figure}[H]
    \centering
    \fbox{\parbox{0.9\textwidth}{\centering\vspace{5cm}[Image de la page compte utilisateur]\vspace{5cm}}}
    \caption{Page compte utilisateur avec onglets commandes et adresses}
    \label{fig:account}
\end{figure}

\subsection{Page de Recherche}

La page de recherche (\texttt{/shop/search}) permet aux utilisateurs de rechercher des produits en utilisant des mots-clés.

\subsubsection{Description}

La page de recherche offre une interface de recherche avec suggestions en temps réel, permettant aux utilisateurs de trouver rapidement les produits souhaités. Les résultats de recherche affichent les produits correspondants aux critères de recherche de manière claire et organisée. Des options de filtrage supplémentaires permettent d'affiner les résultats selon différents critères. Lorsqu'aucun résultat n'est trouvé, un message informatif guide l'utilisateur vers une nouvelle recherche ou la navigation dans le catalogue.

\subsubsection{Capture d'Écran}

\begin{figure}[H]
    \centering
    \fbox{\parbox{0.9\textwidth}{\centering\vspace{5cm}[Image de la page de recherche]\vspace{5cm}}}
    \caption{Page de recherche de produits}
    \label{fig:search}
\end{figure}

\subsection{Barre de Navigation Client}

La barre de navigation (Appbar) est présente sur toutes les pages client et offre un accès rapide aux fonctionnalités principales.

\subsubsection{Description}

La barre de navigation présente le logo de la marque "FLINT" avec un style moderne et cliquable pour retourner à l'accueil. Le menu de navigation offre des liens vers les principales sections : Accueil, Collection, Hommes, Femmes, Enfants, Accessoires, et Chaussures, permettant une navigation intuitive dans toute la plateforme. Une barre de recherche avec icône permet une recherche rapide accessible depuis n'importe quelle page.

Les icônes d'action incluent une icône Panier avec compteur d'articles en temps réel, et un menu utilisateur avec dropdown contenant les liens vers le Profil (compte), les Commandes, et la Déconnexion. Sur les écrans mobiles, un menu hamburger offre une navigation complète adaptée aux petits écrans.

\subsubsection{Fonctionnalités}

La barre de navigation offre une navigation fluide entre les sections, facilitant le parcours utilisateur. La recherche rapide est accessible depuis n'importe quelle page, permettant de trouver rapidement des produits. L'accès rapide au panier avec compteur en temps réel informe l'utilisateur du nombre d'articles dans son panier. Le menu utilisateur regroupe toutes les options de compte en un seul endroit. Le design responsive avec adaptation mobile garantit une expérience optimale sur tous les appareils. L'effet de transparence avec backdrop blur crée un design moderne et élégant.

\subsubsection{Capture d'Écran}

\begin{figure}[H]
    \centering
    \fbox{\parbox{0.9\textwidth}{\centering\vspace{5cm}[Image de la barre de navigation]\vspace{5cm}}}
    \caption{Barre de navigation client}
    \label{fig:navbar}
\end{figure}

\subsection{Pied de Page Client}

Le pied de page (Footer) fournit des informations importantes et des liens utiles.

\subsubsection{Description}

Le pied de page fournit des informations importantes et des liens utiles pour compléter l'expérience utilisateur. La section Logo et Description présente la marque avec une description, renforçant l'identité de la plateforme. Les liens rapides offrent une navigation vers les pages importantes, facilitant l'accès aux sections clés. Les informations de contact incluent l'adresse, le téléphone et l'email, permettant aux clients de contacter facilement l'entreprise.

Les réseaux sociaux sont représentés par des icônes vers les plateformes sociales avec effets de survol, encourageant l'engagement sur les réseaux. Les liens légaux pointent vers les conditions d'utilisation et la politique de confidentialité, assurant la transparence et la conformité. L'information de copyright protège les droits de la plateforme.

\subsubsection{Design}

Le design du pied de page utilise un fond avec dégradé sombre moderne, créant un contraste élégant avec le reste de l'interface. L'organisation en colonnes responsive s'adapte parfaitement à tous les écrans. Les effets de survol sur les liens et icônes améliorent l'interactivité. Le design est cohérent avec le reste de l'interface, assurant une expérience visuelle harmonieuse.

\subsubsection{Capture d'Écran}

\begin{figure}[H]
    \centering
    \fbox{\parbox{0.9\textwidth}{\centering\vspace{5cm}[Image du pied de page]\vspace{5cm}}}
    \caption{Pied de page client}
    \label{fig:footer}
\end{figure}

\newpage

\section{Interfaces Administrateur}

Cette section décrit toutes les interfaces réservées aux administrateurs de la plateforme. Ces interfaces permettent une gestion complète de la boutique en ligne.

\subsection{Tableau de Bord Administrateur}

Le tableau de bord (\texttt{/admin/dashboard}) est la page principale de l'interface d'administration, offrant une vue d'ensemble complète de l'activité de la plateforme.

\subsubsection{Description}

Le tableau de bord administrateur offre une vue d'ensemble complète de l'activité de la plateforme. L'en-tête présente le titre "Tableau de Bord" avec une description, accompagné d'un sélecteur de période permettant de choisir entre 7 jours, 30 jours, 90 jours, ou une période personnalisée pour analyser les données selon les besoins.

Les onglets de navigation organisent les informations en quatre sections principales : Vue d'ensemble (Overview), Revenus (Revenue), Commandes (Orders), et Produits (Products), permettant une navigation claire entre les différents aspects de l'activité.

Les cartes de statistiques affichent les métriques clés dans des cartes modernes : le Revenu Total avec icône et variation en pourcentage, le Nombre de Commandes avec tendance, les Produits Vendus avec statistiques, les Clients Actifs, le Montant des Réductions, et le Taux de Conversion. Ces informations sont présentées de manière visuelle et facilement compréhensible.

Les graphiques offrent une visualisation des données sur la période sélectionnée : graphique de revenus, graphique de commandes par jour, graphique de performance des produits, et graphique de revenus par catégorie. Les tableaux de données complètent l'analyse avec le top des produits vendus, les revenus par produit, et les catégories les plus performantes.

\subsubsection{Fonctionnalités}

Le tableau de bord offre un filtrage par période avec mise à jour dynamique des données, permettant une analyse en temps réel de l'activité. La visualisation graphique des tendances facilite la compréhension des évolutions. Le calcul automatique des statistiques et métriques élimine les erreurs de calcul manuel. Le design moderne avec cartes et graphiques interactifs crée une interface professionnelle et agréable à utiliser. Une fonctionnalité d'export des données est prévue pour l'avenir, permettant une analyse approfondie externe.

\subsubsection{Capture d'Écran}

\begin{figure}[H]
    \centering
    \fbox{\parbox{0.9\textwidth}{\centering\vspace{5cm}[Image du tableau de bord administrateur]\vspace{5cm}}}
    \caption{Tableau de bord administrateur avec statistiques et graphiques}
    \label{fig:admin-dashboard}
\end{figure}

\subsection{Page Gestion des Produits}

La page de gestion des produits (\texttt{/admin/products}) permet aux administrateurs de gérer l'ensemble du catalogue de produits.

\subsubsection{Description}

La page de gestion des produits permet aux administrateurs de gérer l'ensemble du catalogue de produits de manière efficace. L'en-tête présente le titre "Gestion des Produits" avec une description, accompagné d'un bouton "Ajouter un Produit" avec style avec dégradé moderne, facilitant l'ajout rapide de nouveaux articles.

Les badges de filtres actifs affichent les filtres actuellement appliqués avec possibilité de les supprimer, permettant un contrôle précis de l'affichage. Le tableau des produits présente toutes les informations essentielles dans un format moderne : l'image du produit en rond avec bordure, le nom du produit, la catégorie et marque avec badges colorés, le prix avec prix de vente si en promotion, le stock disponible, des badges "Vedette" et "Nouveau" pour identifier les produits spéciaux, et des actions d'édition et suppression.

Le formulaire d'ajout et de modification s'ouvre dans un dialogue modal avec un formulaire complet incluant les informations de base (nom, description, catégorie, marque), la gestion des prix (prix original, prix de vente, réduction), la gestion du stock par taille, l'upload d'images multiples via Cloudinary, l'option "Produit vedette", et la validation des champs. Le dialogue de confirmation de suppression présente un aperçu du produit à supprimer, un message de confirmation, et les boutons d'action pour éviter les suppressions accidentelles. Un état vide avec message et call-to-action s'affiche si aucun produit n'est présent.

\subsubsection{Fonctionnalités}

La page de gestion des produits offre un CRUD complet (Créer, Lire, Modifier, Supprimer) permettant une gestion totale du catalogue. Le filtrage et la recherche de produits facilitent la localisation d'articles spécifiques. L'upload et la gestion d'images via Cloudinary assurent un stockage fiable et performant. La gestion du stock par taille permet un contrôle précis de l'inventaire. La mise en avant de produits permet de promouvoir les articles importants. La validation complète des formulaires prévient les erreurs de saisie. Le design moderne avec tableaux et dialogues crée une interface professionnelle et agréable à utiliser.

\subsubsection{Capture d'Écran}

\begin{figure}[H]
    \centering
    \fbox{\parbox{0.9\textwidth}{\centering\vspace{5cm}[Image de la page gestion des produits]\vspace{5cm}}}
    \caption{Page de gestion des produits avec tableau et formulaire}
    \label{fig:admin-products}
\end{figure}

\subsection{Page Gestion des Commandes}

La page de gestion des commandes (\texttt{/admin/orders}) permet aux administrateurs de visualiser et gérer toutes les commandes passées sur la plateforme.

\subsubsection{Description}

La page de gestion des commandes permet aux administrateurs de visualiser et gérer toutes les commandes passées sur la plateforme. L'en-tête présente le titre "Gestion des Commandes" avec une description, accompagné d'un filtre par statut avec dropdown pour affiner l'affichage selon les besoins.

Le tableau des commandes affiche toutes les informations essentielles dans un format moderne : le numéro de commande formaté et cliquable, les informations du client (nom et email), la date de commande formatée, les images empilées des produits commandés, le montant total, le statut avec badge coloré (En attente en jaune, Confirmée en bleu, Expédiée en violet, Livrée en vert, Annulée en rouge), et les actions pour voir les détails et modifier le statut.

Le dialogue de détails de commande offre une vue complète avec les informations complètes de la commande, la liste détaillée des produits avec quantités, l'adresse de livraison, la méthode de paiement, l'historique des changements de statut, et la possibilité de modifier le statut directement depuis le dialogue. Un état vide avec message informatif s'affiche si aucune commande n'est présente.

\subsubsection{Fonctionnalités}

La page de gestion des commandes offre la visualisation de toutes les commandes en un seul endroit, facilitant le suivi global. Le filtrage par statut permet de se concentrer sur les commandes nécessitant une attention particulière. Les détails complets de chaque commande sont accessibles via un dialogue modal, offrant toutes les informations nécessaires sans quitter la page principale. La modification du statut des commandes peut être effectuée directement depuis l'interface, permettant une mise à jour rapide. Le suivi de l'historique des commandes offre une traçabilité complète. Le design moderne avec badges et tableaux crée une interface professionnelle et facile à utiliser.

\subsubsection{Capture d'Écran}

\begin{figure}[H]
    \centering
    \fbox{\parbox{0.9\textwidth}{\centering\vspace{5cm}[Image de la page gestion des commandes]\vspace{5cm}}}
    \caption{Page de gestion des commandes avec tableau et filtres}
    \label{fig:admin-orders}
\end{figure}

\subsection{Page Gestion des Bannières}

La page de gestion des bannières (\texttt{/admin/banners}) permet aux administrateurs de gérer les bannières promotionnelles affichées sur la page d'accueil.

\subsubsection{Description}

La page de gestion des bannières permet aux administrateurs de gérer les bannières promotionnelles affichées sur la page d'accueil. L'en-tête présente le titre "Gestion des Bannières" avec une description et des instructions d'utilisation pour guider les administrateurs.

La section d'upload offre une carte avec zone de dépôt de fichiers, un bouton de sélection de fichier, une prévisualisation de l'image avant upload, et une validation du format et de la taille pour assurer la qualité des bannières. La grille des bannières affiche toutes les bannières actives avec l'image de la bannière et un effet de survol, un bouton de suppression avec overlay au survol, et une confirmation avant suppression pour éviter les suppressions accidentelles.

Un état vide avec message informatif et icône, accompagné d'un call-to-action pour ajouter la première bannière, guide les administrateurs lors de la première utilisation.

\subsubsection{Fonctionnalités}

La page de gestion des bannières offre l'upload de bannières via Cloudinary, assurant un stockage fiable et performant. La prévisualisation avant upload permet de vérifier l'image avant publication. La suppression de bannières avec confirmation prévient les suppressions accidentelles. La validation des formats d'image garantit la compatibilité et la qualité. L'affichage en grille responsive s'adapte à tous les écrans. Le design moderne avec cartes et effets de survol crée une interface agréable et professionnelle.

\subsubsection{Capture d'Écran}

\begin{figure}[H]
    \centering
    \fbox{\parbox{0.9\textwidth}{\centering\vspace{5cm}[Image de la page gestion des bannières]\vspace{5cm}}}
    \caption{Page de gestion des bannières avec upload et grille}
    \label{fig:admin-banners}
\end{figure}

\subsection{Barre de Navigation Administrateur}

La barre de navigation de l'interface administrateur offre un accès rapide aux fonctionnalités d'administration.

\subsubsection{Description}

La barre de navigation de l'interface administrateur offre un accès rapide aux fonctionnalités d'administration. Le logo de la plateforme avec indication "Administration" renforce l'identité de l'espace admin. Une barre de recherche globale permet de rechercher dans toute l'interface admin. Les icônes d'action incluent une icône de notifications avec badge de compteur pour les alertes importantes, une icône de paramètres, et un menu utilisateur avec dropdown contenant le profil administrateur, les paramètres, et la déconnexion. Le design utilise un fond avec effet de transparence et backdrop blur pour créer un style moderne et élégant.

\subsubsection{Capture d'Écran}

\begin{figure}[H]
    \centering
    \fbox{\parbox{0.9\textwidth}{\centering\vspace{5cm}[Image de la barre de navigation admin]\vspace{5cm}}}
    \caption{Barre de navigation administrateur}
    \label{fig:admin-navbar}
\end{figure}

\subsection{Sidebar / Menu Latéral Administrateur}

Le menu latéral (Drawer) fournit la navigation principale dans l'interface d'administration.

\subsubsection{Description}

Le menu latéral fournit la navigation principale dans l'interface d'administration. La section Logo présente le logo avec dégradé moderne et le titre "FLINT ADMIN", créant une identité visuelle forte. Le menu de navigation offre des liens vers les sections principales : Tableau de bord, Produits, Commandes, et Bannières, chacun accompagné d'une icône pour une identification rapide.

Les états visuels incluent un état actif avec fond coloré et bordure pour indiquer clairement la page courante, des effets de survol avec transitions fluides pour améliorer l'interactivité, et des icônes pour chaque section facilitant la navigation. La version mobile utilise un drawer rétractable avec fond sombre, une animation d'ouverture et fermeture fluide, et un overlay pour fermer le drawer, garantissant une expérience optimale sur mobile.

\subsubsection{Fonctionnalités}

Le menu latéral offre une navigation claire entre les sections, facilitant l'accès aux différentes fonctionnalités. L'indication visuelle de la page active permet à l'administrateur de savoir immédiatement où il se trouve. Le design responsive avec adaptation mobile garantit une expérience optimale sur tous les appareils. Les animations fluides créent une expérience utilisateur agréable. Le design moderne avec dégradés renforce l'identité visuelle de l'interface.

\subsubsection{Capture d'Écran}

\begin{figure}[H]
    \centering
    \fbox{\parbox{0.9\textwidth}{\centering\vspace{5cm}[Image du menu latéral admin]\vspace{5cm}}}
    \caption{Menu latéral administrateur}
    \label{fig:admin-sidebar}
\end{figure}

\newpage

\section{Interfaces d'Authentification}

Cette section décrit les interfaces dédiées à l'authentification des utilisateurs.

\subsection{Page de Connexion}

La page de connexion (\texttt{/auth/login}) permet aux utilisateurs existants de se connecter à leur compte.

\subsubsection{Description}

La page de connexion permet aux utilisateurs existants de se connecter à leur compte de manière sécurisée et intuitive. Le layout avec sidebar présente une sidebar avec dégradé moderne et logo, un message de bienvenue, et un design visuellement attrayant qui crée une première impression positive.

Le formulaire de connexion comprend un champ Email avec validation pour assurer la saisie correcte, un champ Mot de passe avec fonction d'affichage et masquage pour la sécurité et la commodité, une case à cocher "Se souvenir de moi" pour faciliter les connexions futures, un lien "Mot de passe oublié ?" pour la récupération de compte, et un bouton de connexion avec style moderne. Un message avec lien permet de créer un compte pour les nouveaux utilisateurs. La gestion des erreurs affiche les messages d'erreur de manière claire pour guider l'utilisateur.

\subsubsection{Fonctionnalités}

La page de connexion offre une validation des champs en temps réel, permettant de corriger les erreurs immédiatement. L'authentification sécurisée utilise des tokens JWT pour garantir la sécurité des sessions. La gestion des erreurs avec messages clairs guide l'utilisateur en cas de problème. La redirection automatique après connexion améliore l'expérience utilisateur en dirigeant vers la page appropriée. Le design moderne et responsive garantit une expérience optimale sur tous les appareils.

\subsubsection{Capture d'Écran}

\begin{figure}[H]
    \centering
    \fbox{\parbox{0.9\textwidth}{\centering\vspace{5cm}[Image de la page de connexion]\vspace{5cm}}}
    \caption{Page de connexion avec formulaire}
    \label{fig:login}
\end{figure}

\subsection{Page d'Inscription}

La page d'inscription (\texttt{/auth/signup}) permet aux nouveaux utilisateurs de créer un compte.

\subsubsection{Description}

La page d'inscription permet aux nouveaux utilisateurs de créer un compte de manière simple et sécurisée. Le layout avec sidebar utilise le même design que la page de connexion, assurant une cohérence visuelle dans le processus d'authentification.

Le formulaire d'inscription comprend un champ Prénom, un champ Nom, un champ Email avec validation pour assurer la saisie correcte, un champ Mot de passe avec indicateur de force pour guider l'utilisateur vers un mot de passe sécurisé, un champ Confirmation de mot de passe avec validation de la correspondance pour éviter les erreurs, une case à cocher pour accepter les conditions d'utilisation, et un bouton d'inscription avec style moderne. Un message avec lien permet aux utilisateurs existants de se connecter. La validation complète de tous les champs avant soumission garantit la qualité des données.

\subsubsection{Fonctionnalités}

La page d'inscription offre une validation complète des champs, assurant la saisie correcte de toutes les informations. La vérification de la force du mot de passe guide l'utilisateur vers un mot de passe sécurisé. La validation de l'email garantit le format correct. La création de compte sécurisée protège les données utilisateur. La connexion automatique après inscription améliore l'expérience utilisateur en évitant une étape supplémentaire. Le design moderne et responsive garantit une expérience optimale sur tous les appareils.

\subsubsection{Capture d'Écran}

\begin{figure}[H]
    \centering
    \fbox{\parbox{0.9\textwidth}{\centering\vspace{5cm}[Image de la page d'inscription]\vspace{5cm}}}
    \caption{Page d'inscription avec formulaire}
    \label{fig:signup}
\end{figure}

\newpage

\section{Pages Supplémentaires}

\subsection{Page 404 - Page Non Trouvée}

La page 404 (\texttt{/*}) s'affiche lorsque l'utilisateur tente d'accéder à une URL inexistante.

\subsubsection{Description}

La page 404 s'affiche lorsque l'utilisateur tente d'accéder à une URL inexistante. Un message clair indique que la page n'existe pas, guidant l'utilisateur vers une solution. Le design moderne est cohérent avec le reste de l'application, assurant une expérience visuelle harmonieuse même en cas d'erreur. Un bouton de retour offre un lien pour retourner à la page d'accueil facilement. Une illustration visuelle rend la page plus attrayante et moins frustrante pour l'utilisateur.

\subsubsection{Capture d'Écran}

\begin{figure}[H]
    \centering
    \fbox{\parbox{0.9\textwidth}{\centering\vspace{5cm}[Image de la page 404]\vspace{5cm}}}
    \caption{Page 404 - Page non trouvée}
    \label{fig:404}
\end{figure}

\subsection{Page Non Autorisée}

La page non autorisée (\texttt{/restrict}) s'affiche lorsque l'utilisateur tente d'accéder à une ressource sans les permissions nécessaires.

\subsubsection{Description}

La page non autorisée s'affiche lorsque l'utilisateur tente d'accéder à une ressource sans les permissions nécessaires. Un message d'avertissement clair explique le problème d'autorisation, informant l'utilisateur de la situation. Une action suggérée offre un lien pour se connecter ou retourner à l'accueil, guidant l'utilisateur vers une solution. Le design est cohérent avec le reste de l'application, assurant une expérience visuelle harmonieuse.

\subsubsection{Capture d'Écran}

\begin{figure}[H]
    \centering
    \fbox{\parbox{0.9\textwidth}{\centering\vspace{5cm}[Image de la page non autorisée]\vspace{5cm}}}
    \caption{Page non autorisée}
    \label{fig:unauthorized}
\end{figure}

\subsection{Pages Légales}

\subsubsection{Conditions d'Utilisation}

La page des conditions d'utilisation (\texttt{/terms-conditions}) présente les termes et conditions d'utilisation de la plateforme.

\subsubsection{Politique de Confidentialité}

La page de politique de confidentialité (\texttt{/privacy-policy}) présente la politique de protection des données personnelles.

\newpage

\section{Conclusion}

\subsection{Résumé du Projet}

Flint Commerce représente une plateforme e-commerce complète et moderne, offrant une expérience utilisateur exceptionnelle tant pour les clients que pour les administrateurs. Le projet a été développé avec une attention particulière portée à la qualité de l'interface utilisateur, à la responsivité, et à la facilité d'utilisation.

L'architecture full-stack, utilisant React.js pour le frontend et Node.js/Express pour le backend, garantit des performances optimales et une maintenabilité élevée. L'intégration de technologies modernes comme Redux Toolkit pour la gestion d'état, Tailwind CSS pour le styling, et MongoDB pour la base de données, assure une base solide pour l'évolution future de la plateforme.

\subsection{Points Forts}

Les principaux points forts de cette plateforme incluent :

\begin{itemize}
    \item \textbf{Interface Utilisateur Moderne} : Design contemporain avec dégradés, animations fluides, et une esthétique visuelle attrayante qui améliore l'expérience utilisateur.
    
    \item \textbf{Responsive Design} : Toutes les interfaces sont entièrement responsive, garantissant une expérience optimale sur desktop, tablette et mobile.
    
    \item \textbf{Fonctionnalités Complètes} : La plateforme couvre tous les aspects d'une solution e-commerce moderne, de la navigation des produits à la gestion administrative complète.
    
    \item \textbf{Sécurité} : Implémentation de l'authentification JWT, validation des données, et gestion sécurisée des transactions.
    
    \item \textbf{Performance} : Architecture optimisée avec gestion d'état efficace, chargement paresseux, et requêtes API optimisées.
    
    \item \textbf{Expérience Utilisateur} : Processus d'achat intuitif, gestion de panier fluide, et interface administrative complète pour la gestion de la boutique.
\end{itemize}

\subsection{Interfaces Développées}

Ce projet comprend un total de plus de 20 interfaces distinctes, organisées en trois espaces principaux :

\begin{enumerate}
    \item \textbf{Espace Client} : 9 interfaces principales (Accueil, Catalogue, Détails Produit, Panier, Paiement, Compte, Recherche, Navigation, Footer)
    
    \item \textbf{Espace Administrateur} : 6 interfaces principales (Tableau de bord, Produits, Commandes, Bannières, Navigation, Sidebar)
    
    \item \textbf{Espace Authentification} : 2 interfaces (Connexion, Inscription)
    
    \item \textbf{Pages Supplémentaires} : Pages d'erreur et pages légales
\end{enumerate}

Chaque interface a été conçue avec soin pour offrir une expérience utilisateur optimale, avec une attention particulière portée à la cohérence visuelle, à la facilité de navigation, et à l'accessibilité.

\subsection{Perspectives d'Évolution}

La plateforme Flint Commerce offre une base solide pour des évolutions futures, notamment :

\begin{itemize}
    \item Intégration de méthodes de paiement supplémentaires
    \item Système de recommandation de produits basé sur l'IA
    \item Application mobile native
    \item Programme de fidélité et de récompenses
    \item Chat en direct pour le support client
    \item Système de reviews et ratings avancé
    \item Multi-langues et multi-devises
    \item Analytics avancés pour les administrateurs
\end{itemize}

\subsection{Conclusion Générale}

En conclusion, Flint Commerce représente une implémentation réussie d'une plateforme e-commerce moderne et complète. Le projet démontre une maîtrise des technologies web modernes, une attention aux détails dans la conception de l'interface utilisateur, et une compréhension approfondie des besoins tant des clients que des administrateurs d'une plateforme e-commerce.

La qualité du code, la structure de l'architecture, et l'expérience utilisateur soignée font de cette plateforme une solution e-commerce robuste et prête pour la production. Le projet constitue une base excellente pour une boutique en ligne professionnelle, avec toutes les fonctionnalités nécessaires pour gérer efficacement les ventes en ligne.

\vspace{1cm}

\begin{center}
\textit{Fin du Rapport}
\end{center}

\end{document}

